\section{A Framework for Open-World Agents}
\label{sec:agent_framework}

Having defined the environment and the nature of novelty, we now turn to the structure of the agent itself. The agents studied in this dissertation operate by combining symbolic reasoning at the level of task planning with subsymbolic execution at the level of environment interaction and execution. 
This section formalizes the components of this framework: symbolic planning, reinforcement learning, and their integration into a hybrid architecture. More specifically, we will provide the preliminary definitions of the components of the agent that will be used throughout the dissertation.
% Rather than presenting these as independent formalisms, we develop them as complementary components of a single agent that must reason, act, and adapt under novelty.

\subsection{Symbolic Planning}
\label{sec:symbolic_planning}

The agent's high-level decision-making is conducted through symbolic planning, which enables structured reasoning over abstract representations of states and actions. Following the representation introduced in \cref{sec:state_representations}, we adopt a STRIPS-style formalism \citep{fikes1971strips} in which symbolic states, preconditions, and effects are defined over grounded literals. This representation is used throughout the dissertation and provides the foundation for the operator--executor consistency conditions developed in \cref{ch:rapid_learn}.

\paragraph{Symbolic states.}
Let $\Language$ be the agent's formal language of predicates, variables, constants, and possibly functions. A \emph{grounded literal} is an atom $p(c_1, \ldots, c_k)$ or its negation $\neg p(c_1, \ldots, c_k)$, where $p$ is a predicate symbol and each $c_i$ is a constant. We write $\Lit$ for the set of all grounded literals over $\Language$. A symbolic state $\SymState \in \SymStates$ is a consistent subset $\SymState \subseteq \Lit$ containing no atom together with its negation. Under the closed-world convention (\cref{sec:state_representations}), literals not contained in $\SymState$ are treated as false. In propositional planning domains, $\SymState$ is typically represented by the set of positive ground atoms that hold; negated literals are implicit. A conjunction of positive literals $\Phi$ is satisfied in $\SymState$, written $\SymState \models \Phi$, if and only if $\Phi \subseteq \SymState$.

\paragraph{Planning domain and task.}
The agent's symbolic model is organized into a \emph{planning domain} and a \emph{planning task}. The planning domain $\PlanningDomain = \langle \Fluents, \Operators \rangle$ specifies the domain vocabulary: $\Fluents$ is the set of fluent symbols (predicate and function symbols whose grounded instances may change over time), and $\Operators$ is a set of operator schemas. A planning task is a tuple $\PlanningTask = \langle \PlanningDomain, \SymState_0, \Goal \rangle$, equivalently written $\PlanningTask = \langle \Fluents, \Operators, \SymState_0, \Goal \rangle$, where $\SymState_0 \in \SymStates$ is the initial state and $\Goal$ is a goal condition over literals (a conjunction of positive literals, or a disjunctive normal form thereof). Throughout the dissertation, we write $\Pre_\Operator = \psi_\Operator$ for the precondition literal set of operator $\Operator$ and $\Eff_\Operator = \omega_\Operator$ for its effect literal set.

\paragraph{Operators.}
Each operator schema defines a parameterized action template. A \emph{grounded operator} $\Operator \in \Operators$ is specified as
\[
\Operator = \langle \Pre_\Operator, \Eff_\Operator^+, \Eff_\Operator^- \rangle,
\]
where $\Pre_\Operator \subseteq \Lit$ is a set of positive literals whose conjunction must hold for the operator to be applicable, $\Eff_\Operator^+ \subseteq \Lit$ is the \emph{add list} of positive literals made true by the operator, and $\Eff_\Operator^- \subseteq \Lit$ is the \emph{delete list} of positive literals made false. When an operator's effects are referenced as a single set---for example, in effect satisfaction checks during execution---$\Eff_\Operator$ denotes the postcondition add list $\Eff_\Operator^+$. Operator schemas are instantiated over domain objects to produce grounded operators; for instance, a schema \texttt{break(?x)} grounded with a tree constant yields a grounded operator with concrete preconditions and effects. Domains with numeric fluents extend this representation with assignment and arithmetic effects, as in PDDL~\citep{mcdermott1998pddl}; the domains evaluated in this dissertation are predominantly propositional.

An operator $\Operator$ is \emph{applicable} in symbolic state $\SymState$ if $\SymState \models \Pre_\Operator$, equivalently $\Pre_\Operator \subseteq \SymState$. The operator's effects are \emph{satisfied} in $\SymState$, written $\SymState \models \Eff_\Operator$, if $\Eff_\Operator^+ \subseteq \SymState$ and $\Eff_\Operator^- \cap \SymState = \emptyset$. A \emph{partial effect failure} occurs when a strict subset of $\Eff_\Operator^+$ is observed after execution; a \emph{total effect failure} occurs when none of $\Eff_\Operator^+$ is observed.

\paragraph{State transitions and plans.}
Applying an operator induces a deterministic state transition. The \emph{application function} $\Apply : \SymStates \times \Operators \rightarrow \SymStates$ is defined by
\[
\Apply(\SymState, \Operator) =
\begin{cases}
(\SymState \setminus \Eff_\Operator^-) \cup \Eff_\Operator^+ & \text{if } \Pre_\Operator \subseteq \SymState, \\
\text{undefined} & \text{otherwise}.
\end{cases}
\]
A \emph{plan} is a finite sequence of grounded operators $\Plan = [\Operator_1, \dots, \Operator_n]$. Executing $\Plan$ from $\SymState_0$ produces the sequence of states $\SymState_{i+1} = \Apply(\SymState_i, \Operator_i)$ for $i = 1, \ldots, n$, provided each operator is applicable in its respective state. The plan is \emph{valid} (or a \emph{planning solution}) if every operator is applicable and the final state satisfies the goal: $\SymState_n \models \Goal$.

\paragraph{Operator dependencies.}
Operators within a plan may depend on one another through shared literals. Operator $\Operator_j$ \emph{depends on} operator $\Operator_i$, written $\Operator_i \Dep \Operator_j$, if the effects of $\Operator_i$ are contained in the preconditions of $\Operator_j$:
\[
\Operator_i \Dep \Operator_j \iff \Eff_{\Operator_i} \subseteq \Pre_{\Operator_j}.
\]
This relation is used in \cref{ch:rapid_learn} to identify downstream operators affected by a failed execution and to define \emph{plannable states} from which planning can resume.

Classical planning assumes deterministic transitions and complete, correct domain models. Planning domains and problems are commonly specified using the Planning Domain Definition Language (PDDL) \citep{mcdermott1998pddl}, and solved using domain-independent heuristic search algorithms such as Fast Downward \citep{helmert2006fast} or Metric-FF \citep{ecai02}.

In the context of this dissertation, symbolic planning provides the agent with the ability to decompose long-horizon tasks into sequences of abstract actions. This decomposition is valuable precisely because it creates structure that can be exploited when novelty occurs: if an operator fails, the agent can identify \emph{which} step of the plan was affected and reason about what must change.

\subsection{Reinforcement Learning}
\label{sec:rl}

While symbolic planning reasons over abstract state descriptions, the agent must ultimately execute actions in the environment, where internal state representations may be high-dimensional, observations partial, and dynamics stochastic. Reinforcement learning provides the framework for learning effective behaviors through direct interaction~\citep{sutton2018reinforcement}. At the execution level, $\States$ denotes the agent's internal state representation---which may consist of sensor readings, learned features, or structured numerical descriptions---rather than the full environment state $e \in \EnvStates$ (\cref{sec:state_representations}).

\paragraph{Markov decision processes.}
The agent's interaction with the environment at the primitive level is modeled as a Markov Decision Process (MDP), defined as $M = \langle \States, \Actions, \Transition, \Reward, \Discount \rangle$, where $\States$ is the set of states, $\Actions$ is the set of primitive actions, $\Transition(s' \mid s, a)$ defines the probability of transitioning to state $s'$ after taking action $a$ in state $s$, $\Reward(s, a, s')$ is the immediate reward received upon this transition, and $\Discount \in [0, 1)$ is the discount factor. The transition dynamics and reward satisfy the \emph{Markov property}: the distribution over next states and rewards depends only on the current state and action, not on prior history. When transitions are deterministic, $\Transition(s' \mid s, a) = 1$ for a unique next state $s'$ and $0$ otherwise; deterministic transition functions, as used in gridworld domains (\cref{ch:benchmarks}), are a special case of this formulation. At each time step $t$, the agent observes state $s_t \in \States$, selects action $a_t \in \Actions$, transitions to $s_{t+1}$ according to $\Transition(\cdot \mid s_t, a_t)$, and receives reward $r_{t+1} = \Reward(s_t, a_t, s_{t+1})$.

\paragraph{Policies.}
A policy $\Policy(a \mid s)$ specifies the agent's behavior by defining a probability distribution over actions given a state; deterministic policies are the special case where $\Policy(a \mid s) \in \{0, 1\}$ for each $(s, a)$ pair. The agent's objective is to find a policy that maximizes expected cumulative discounted reward. The \emph{return} from time step $t$ is $G_t = \sum_{k=0}^{\infty} \Discount^k r_{t+k+1}$, the performance objective is $J(\Policy) = \Expect_\Policy[G_0 \mid s_0]$, and an \emph{optimal policy} satisfies $\Policy^* = \arg\max_\Policy J(\Policy)$. The state-value function is $\Value^\Policy(s) = \Expect_\Policy[G_t \mid s_t = s]$, and the action-value function is $\QValue^\Policy(s,a) = \Expect_\Policy[G_t \mid s_t = s, a_t = a]$. These satisfy the Bellman expectation equations:
\begin{align}
\Value^\Policy(s) &= \Expect_\Policy\left[r_{t+1} + \Discount \Value^\Policy(s_{t+1}) \mid s_t = s\right], \\
\QValue^\Policy(s,a) &= \Expect_\Policy\left[r_{t+1} + \Discount \Expect_{a' \sim \Policy(\cdot \mid s_{t+1})} \QValue^\Policy(s_{t+1}, a') \mid s_t = s, a_t = a\right].
\end{align}
In discrete settings with finite $\States$ and $\Actions$, value functions can be represented in tabular form; in continuous settings where $\States \subseteq \Real^n$ and $\Actions \subseteq \Real^m$, policies and value functions are approximated using parameterized function approximators $\Policy_\theta(a \mid s)$, $\Value_\phi(s)$, or $\QValue_\phi(s, a)$, optimized by algorithms such as DQN~\citep{mnih2015human}, PPO~\citep{schulman2017proximal}, and SAC~\citep{haarnoja2018soft}.

\paragraph{Semi-Markov decision processes and temporal abstraction.}
Long-horizon tasks require decisions at multiple levels of temporal granularity. \citet{sutton1999options} formalize temporal abstraction in the framework introduced in \emph{Between MDPs and Semi-MDPs}: given a base MDP $M = \langle \States, \Actions, \Transition, \Reward, \Discount \rangle$, an \emph{option} $\omega = \langle \mathcal{I}_\omega, \Policy_\omega, \beta_\omega \rangle$ consists of an initiation set $\mathcal{I}_\omega \subseteq \States$ specifying the states in which the option may be initiated, an intra-option policy $\Policy_\omega(a \mid s)$ that selects primitive actions while the option executes, and a termination function $\beta_\omega(s) \in [0, 1]$ giving the probability that the option terminates upon entering state $s$. When the agent selects an option $\omega$ in a state $s \in \mathcal{I}_\omega$, the intra-option policy executes primitive actions; after each transition to $s'$, the option terminates with probability $\beta_\omega(s')$ or continues. A set of options $\Omega$ together with a policy over options $\mu(\omega \mid s)$ defines a \emph{semi-Markov decision process} (SMDP) over $\States$, in which high-level decisions select options and each option induces a semi-Markov transition of random duration~\citep{sutton1999options,sutton2018reinforcement}. Primitive actions are the degenerate special case of options that are available in every state, follow a single primitive action, and terminate with probability one.

From the agent's perspective, an option behaves like a single macro-action: instead of selecting a primitive action at each time step, the agent selects an option and delegates the intervening steps to the option's policy. For example, rather than planning over individual \texttt{move-forward} and \texttt{turn-left} actions, the agent may select a \texttt{navigate-to-tree} option that internally executes many primitive steps before terminating. \Cref{sec:hybrid} instantiates this formalism through \emph{executors}---options tied to symbolic operators---and couples the resulting SMDP to a symbolic planning task.
\subsection{Hybrid Planning and Learning}
\label{sec:hybrid}

The symbolic planning formalism (\cref{sec:symbolic_planning}) and the reinforcement learning formalism (\cref{sec:rl}) operate at different levels of abstraction. Symbolic planning reasons over discrete operators and fluents to produce long-horizon plans; reinforcement learning operates over the agent's subsymbolic state representation to select primitive actions or temporally extended options. A hybrid architecture integrates both: the agent plans over symbolic operators at the high level and executes those operators through an SMDP at the low level, where each operator is realized by an executor.

In a hybrid system, decision-making operates over two coupled representations. At the symbolic level, the agent uses its planning task $\PlanningTask = \langle \PlanningDomain, \SymState_0, \Goal \rangle$ to generate a plan $\Plan = [\Operator_1, \dots, \Operator_n]$. At the execution level, a base MDP $M = \langle \States, \Actions, \Transition, \Reward, \Discount \rangle$ models primitive interaction with the environment, and a set of executors $\Executors$ defines the temporally extended actions of the induced SMDP. An abstraction function $\Abstraction : \States \rightarrow \SymStates$ maps subsymbolic states to symbolic states represented as sets of grounded literals (\cref{sec:symbolic_planning}), connecting the two levels of representation.

The critical link between planning and execution is the \emph{executor}. Each operator $\Operator \in \Operators$ is associated with an executor $\Executor_\Operator \in \Executors$ that realizes the operator's intended effects in the environment. An executor is an option $\Executor_\Operator = \langle I_\Operator, \Policy_\Operator, \beta_\Operator \rangle$, where $I_\Operator \subseteq \States$ is the initiation set, $\Policy_\Operator(a \mid s)$ is the intra-option policy, and $\beta_\Operator(s_0, s) \in \{0, 1\}$ is a termination indicator specifying whether the executor may terminate in state $s$ when initiated in $s_0$. The additional constraint beyond the standard options framework is that each executor is tied to a symbolic operator whose preconditions and effects provide its semantic meaning~\citep{goel2022rapidlearn}.

A symbolic plan $\Plan = [\Operator_1, \dots, \Operator_n]$ thus induces a sequence of executors $[\Executor_{\Operator_1}, \dots, \Executor_{\Operator_n}]$ that are carried out sequentially in the SMDP. The hybrid system is defined as $\Hybrid = \langle \PlanningTask, M, \Abstraction, \Executors \rangle$, formalizing the coupling of symbolic planning over operators with semi-Markov execution through learned and hand-coded policies. \Cref{fig:agent_environment} illustrates this two-level architecture.

% Agent-environment interaction figure
\begin{figure}[t]
    \centering
    \begin{tikzpicture}[
        >=Stealth,
        % --- Node styles (Ch.4 palette) ---
        symstate/.style={
            circle, draw=blue!60!black, fill=blue!8,
            minimum size=0.95cm, font=\small, thick
        },
        envstate/.style={
            circle, draw=teal!70!black, fill=teal!8,
            minimum size=0.95cm, font=\small, thick
        },
        opnode/.style={
            rectangle, rounded corners=3pt, draw=blue!60!black, fill=blue!8,
            minimum width=1.3cm, minimum height=0.65cm,
            font=\small, align=center, thick
        },
        exnode/.style={
            rectangle, rounded corners=3pt, draw=teal!70!black, fill=teal!8,
            minimum width=1.3cm, minimum height=0.65cm,
            font=\small, align=center, thick
        },
        goalnode/.style={
            rectangle, rounded corners=3pt, draw=green!60!black, fill=green!8,
            minimum width=1.1cm, minimum height=0.65cm,
            font=\small, align=center, thick
        },
        % --- Arrow styles ---
        planflow/.style={->, thick, blue!50!black},
        exflow/.style={->, thick, teal!50!black},
        mapline/.style={<->, densely dashed, gray!60!black, thin},
        corr/.style={<->, densely dashed, orange!70!black, semithick},
        % --- Labels ---
        levellabel/.style={font=\footnotesize\itshape, gray!60!black},
        alphabox/.style={font=\footnotesize, gray!60!black},
        % --- Legend ---
        legendbox/.style={minimum width=0.45cm, minimum height=0.3cm, rounded corners=2pt, thick, inner sep=0pt},
        legendlabel/.style={font=\scriptsize, anchor=west, inner sep=0pt},
    ]

    % ===== SYMBOLIC LEVEL (top) =====
    \node[levellabel, anchor=west] at (-5.0, 0.2) {Symbolic};
    \draw[gray!40, thin, densely dashed] (-5.0, 0.0) -- (5.6, 0.0);

    % Symbolic states
    \node[symstate] (s0) at (-3.2, -0.9) {$\SymState_0$};
    \node[symstate] (s1) at (-0.2, -0.9) {$\SymState_1$};
    \node[symstate] (s2) at (2.8, -0.9) {$\SymState_2$};

    % Operators (positioned on the arrows)
    \draw[planflow] (s0) -- (s1) node[midway, above=3pt, opnode] (o1) {$\Operator_1$};
    \draw[planflow] (s1) -- (s2) node[midway, above=3pt, opnode] (o2) {$\Operator_2$};

    % Plan brace
    \draw[decorate, decoration={brace, amplitude=4pt, raise=2pt}, blue!40]
        (o1.north west) -- (o2.north east)
        node[midway, above=7pt, font=\small, blue!60!black] {Plan $\Plan = [\Operator_1, \Operator_2]$};

    % ===== ABSTRACTION MAPPING (middle) =====
    \node[alphabox] at (5.2, -2.0) {$\Abstraction$};
    \draw[<-, densely dashed, gray!60!black, semithick] (4.9, -1.2) -- (4.9, -2.65);

    % ===== SUBSYMBOLIC LEVEL (bottom) =====
    \node[levellabel, anchor=west] at (-5.0, -2.35) {Subsymbolic};
    \draw[gray!40, thin, densely dashed] (-5.0, -2.55) -- (5.6, -2.55);

    % Environment states
    \node[envstate] (e0) at (-3.2, -3.45) {$e_0$};
    \node[envstate] (e1) at (-0.2, -3.45) {$e_1$};
    \node[envstate] (e2) at (2.8, -3.45) {$e_2$};

    % Executors (positioned on the arrows)
    \draw[exflow] (e0) -- (e1) node[midway, below=3pt, exnode] (x1) {$\Executor_{\Operator_1}$};
    \draw[exflow] (e1) -- (e2) node[midway, below=3pt, exnode] (x2) {$\Executor_{\Operator_2}$};

    % ===== CORRESPONDENCE ARROWS =====
    \draw[corr] (o1.south west) -- (x1.north west)
        node[midway, left=2pt, font=\scriptsize, orange!70!black] {};
    \draw[corr] (o2.south east) -- (x2.north east)
        node[midway, right=2pt, font=\scriptsize, orange!70!black] {};

    % ===== LEGEND (centered, evenly spaced) =====
    \draw[gray!30, thin] (-5.0, -4.75) -- (5.6, -4.75);

    \node[legendbox, draw=blue!60!black, fill=blue!8]    (L1) at (-3.8, -5.1) {};
    \node[legendlabel] at (-3.5, -5.1) {Symbolic};

    \node[legendbox, draw=teal!70!black, fill=teal!8]     (L2) at (-1.0, -5.1) {};
    \node[legendlabel] at (-0.7, -5.1) {Subsymbolic};

    \draw[corr] (1.6, -5.1) -- (2.1, -5.1);
    \node[legendlabel] at (2.2, -5.1) {Operator--executor pair};

    \end{tikzpicture}
    \caption[Agent--environment interaction in the hybrid framework]{The two-level structure of the hybrid planning and learning framework. \textbf{Top (symbolic):} The agent generates a plan $\Plan$ as a sequence of operators $\Operator_i$ transitioning between symbolic states $\SymState$. \textbf{Bottom (subsymbolic):} The environment consists of concrete states $e$, and transitions are realized by executors $\Executor_{\Operator_i}$---policies that implement the corresponding operators. The abstraction function $\Abstraction : \States \to \SymStates$ maps environment states to symbolic descriptions. Each operator--executor pair (dashed orange) links a symbolic action to its realization in the environment. Novelty manifests when the executor for an operator no longer achieves the intended symbolic transition.}
    \label{fig:agent_environment}
\end{figure}

Returning to the crafting example from \cref{sec:running_example}, the agent's plan to craft a pogostick might be $\Plan = [\texttt{break\_tree}, \texttt{craft\_planks}, \texttt{craft\_sticks}, \texttt{craft\_pogostick}]$. Each operator in this plan is associated with an executor: the \texttt{approach\_tree} executor navigates to a tree (can be implemented using a navigation policy) and \texttt{break\_tree} performs the breaking action, while the \texttt{craft\_planks} executor navigates to a crafting table and executes the crafting recipe. The symbolic planner reasons about the sequence; the executors handle the details of carrying it out.

\subsection{Novelty in Hybrid Systems}
\label{sec:novelty_hybrid}

The hybrid framework makes explicit where novelty can manifest and how it propagates through the agent's decision-making.

At the symbolic level, novelty may render the planning model incomplete: new fluents, new operators, or changed preconditions and effects may mean that no valid plan exists within the agent's current model. This corresponds to prohibitive novelty (\cref{def:prohibitive}).

At the execution level, novelty may cause an executor to fail: the environment dynamics have changed such that the policy $\Policy_\Operator$ associated with operator $\Operator$ no longer achieves the intended effects, even though the symbolic preconditions $\Pre_\Operator$ are satisfied. This corresponds to obstructive novelty (\cref{def:obstructive}), and it manifests as a discrepancy between what the symbolic model predicts and what actually occurs when actions are executed.

This two-level structure has an important consequence: in a hybrid system, the point of failure is often identifiable. When an executor fails, the agent knows \emph{which} operator was being executed and can examine the discrepancy between expected and observed symbolic states. This localization of failure---enabled by the structured decomposition of behavior into operators and executors---is what makes targeted adaptation possible. Without this structure, a failure in the environment would be a global signal (reward decreased) with no indication of what specifically went wrong.

This observation motivates the central claim of the dissertation: that structured abstractions over actions and interactions determine whether novelty-induced mismatch can be resolved efficiently and whether learned behaviors can generalize across variations. The next section develops this argument formally.